\documentclass{article}[12pt, a4paper]
\usepackage{amsmath}
\usepackage{esint} % Required for \oiint and \oiiint
\usepackage{hyperref}
\newcommand{\gfrac}[2]{\frac{\displaystyle #1}{\displaystyle #2}}
\begin{document}
% \tableofcontents
% \newpage
\centering \section*{Core: Quantum Mechanics Formulations}

%% Schrodinger Equation
\textbf{Schrodinger Equation}
$$\boxed{\frac{-\hbar^2}{2m} \mathbf{\nabla^2} \psi(\mathbf{r}.t) = i\hbar \frac{\partial}{\partial t} \psi(\mathbf{r}.t)} \textbf{ (TDSE)}$$
$$\boxed{\hat{\mathbf{H}} \psi(\mathbf{r}.t) = i\hbar \frac{\partial}{\partial t} \psi(\mathbf{r}.t)} \textbf{ (TDSE)}$$
$$\boxed{\hat{\mathbf{H}} \psi(\mathbf{r}.t) = E \psi(\mathbf{r}.t)} \textbf{ (TISE)}$$ \\

%% Nature of Psi
\textbf{The Wave Funtion}
$$\boxed{\rho (\mathbf{r}, t) = \psi(\mathbf{r}.t)^\mathbf{*} \psi(\mathbf{r}.t) =  {|\psi(\mathbf{r}.t)|}^2} \textbf{ (Probability Density)}$$
$$\boxed{ \int_{-\infty}^{\infty} \rho (\mathbf{r}, t) \mathrm{d}\mahtbf{\tau} = 1} \textbf{ (Normalization Condition)}$$

\centering \section*{Aplication: One Dimentional Potentials}

%% One dimetional infinity potential well
\textbf{Infinify Potential Well}
$$
\boxed{
V(x) =
\begin{cases}
	0 & \text{if } 0 \le x \le a \\
	\infty & \text{Otherwse}
\end{cases}
} \text{ (Potential)}
$$
$$
% \boxed{
% \begin{aligned}
% \psi(0) = \psi (a)  \\
% \\
% \gfrac{\partial \psi(x)}{\partial x}\big|_{x=0} = \gfrac{\partial \psi(x)}{\partial x}\big|_{x=a}
% \end{aligned}
% } \textbf{ (Bonudary Conditions)}
\boxed{
\psi(0) = \psi (a) \quad : \quad
\gfrac{\partial \psi(x)}{\partial x}\big|_{x=0} = \gfrac{\partial \psi(x)}{\partial x}\big|_{x=a}
} \textbf{ (Bonudary Conditions)}
$$
$$
\boxed{}
$$

%% Appendex A
\section*{Appendex A: Probability}
Say, I room having peoples with different age, and $N(j)$ representing number of people with age $j$,
(Eg: if $3$ people with age $25$, then $N(25) = 3$)

$$\boxed{N = \sum N(j)} \text{ (Total number of peopele in the room)}$$
$$\boxed{P(j) = \gfrac{N(j)}{N}} \text{ (Probabiliy of getting age $j$)}$$
$$\boxed{\sum P(j) = 1} \text{ (Total probability)}$$
$$\boxed{\langle j \rangle = \gfrac{\sum jN(j)}{N} = \sum j P(j)} \text{ (Avarage/mean value of $j$)}$$
$$\boxed{\langle j^2 \rangle = \gfrac{\sum j^2 N(j)}{N} = \sum j^2 P(j)} \text{ (Avarage/mean value of $j^2$)}$$
$$\boxed{\langle f(j) \rangle = \gfrac{\sum f(j) N(j)}{N} = \sum f(j) P(j)} \text{ (Avarage/mean value of $f(j)$)}$$
$$\boxed{\sigma_{f(j)} = \langle \Delta f(j) \rangle = \sqrt{\langle f(j)^2 \rangle - \langle f(j) \rangle}} \text{ (SD/Unceritinity of $f(j)$)}$$

\section*{Appendex B: Most Useful Integral}
$$\boxed{\int^{\infty}_0 x^n e^{\displaystyle -\alpha x^m} dx = \gfrac{1}{m} \alpha^{- \big(\gfrac{n+1}{m}\big)} \Gamma \big(\gfrac{n+1}{m}\big)}$$ \\
$(\alpha > 0, m > 0, n > -1)$
$$\boxed{\Gamma(n + 1) = n\Gamma(n) = n!}$$
$$\boxed{\Gamma\big(\gfrac{1}{2}\big) = \sqrt{\pi}}$$

\end{document}
